\documentclass[11pt]{amsart}

\usepackage[T1]{fontenc}
\usepackage[utf8]{inputenc}
\usepackage{lmodern}
\usepackage{amsmath}
\usepackage{amssymb}
\usepackage{amsfonts}
\usepackage{amsthm}
\usepackage{mathtools}
\usepackage{geometry}
\usepackage[hidelinks]{hyperref}

\geometry{margin=1in}
\numberwithin{equation}{section}
\setcounter{tocdepth}{2}
\setcounter{secnumdepth}{2}

\theoremstyle{plain}
\newtheorem{theorem}{Theorem}[section]
\newtheorem{proposition}[theorem]{Proposition}
\newtheorem{lemma}[theorem]{Lemma}
\newtheorem{corollary}[theorem]{Corollary}

\theoremstyle{definition}
\newtheorem{definition}[theorem]{Definition}

\theoremstyle{remark}
\newtheorem{remark}[theorem]{Remark}

\newcommand{\Torus}{\mathbb T}
\newcommand{\dd}{\,\mathrm d}

\title[Prime-Pair Kernels]{Prime-Pair Kernels, Mangoldt Resonances, and a Finite-Height Energy Bridge}
\author{}
\date{}

\begin{document}

\maketitle
\tableofcontents

\section{Shells and normalized kernels}
\label{sec:shells}

Write $\mathbb N_0=\{0,1,2,\ldots\}$.  Let
\begin{equation}
\mathcal S_R\subset\mathbb N_0^2
\qquad(R>0)
\label{eq:shells}
\end{equation}
be finite nested sets satisfying
\begin{equation}
\mathcal S_R\subset\mathcal S_{R'}\quad(R<R'),
\qquad
\bigcup_{R>0}\mathcal S_R=\mathbb N_0^2.
\label{eq:shell_exhaustion}
\end{equation}
The parameter $R$ indexes the arithmetic shell; the letter $T$ will be reserved for the height of zeta zeros.

Fix distinct primes $p,q$.  Define
\begin{align}
W_R^{(p,q)}
&=
\sum_{(a,b)\in\mathcal S_R}p^{-a}q^{-b},
\label{eq:WR}
\\
S_R^{(p,q)}
&=
\sum_{(a,b)\in\mathcal S_R}p^{-2a}q^{-2b},
\label{eq:SR}
\\
w_{a,b;R}^{(p,q)}
&=
\frac{p^{-a}q^{-b}}{W_R^{(p,q)}}
\qquad((a,b)\in\mathcal S_R).
\label{eq:weights}
\end{align}
Then the weights are positive and have sum one.

\begin{proposition}[Coefficient energy]
\label{prop:coefficient_energy}
Set
\begin{equation}
Q_R^{(p,q)}
=
\sum_{(a,b)\in\mathcal S_R}
\left(w_{a,b;R}^{(p,q)}\right)^2
=
\frac{S_R^{(p,q)}}{\left(W_R^{(p,q)}\right)^2}.
\label{eq:QR}
\end{equation}
Then
\begin{equation}
Q_R^{(p,q)}\longrightarrow
Q_\infty^{(p,q)}
=
\frac{p-1}{p+1}\frac{q-1}{q+1}
\in(0,1).
\label{eq:Qinfty}
\end{equation}
\end{proposition}

\begin{proof}
Monotone convergence of the two geometric series gives
\begin{equation}
W_R^{(p,q)}\longrightarrow
\frac{1}{(1-p^{-1})(1-q^{-1})},
\qquad
S_R^{(p,q)}\longrightarrow
\frac{1}{(1-p^{-2})(1-q^{-2})}.
\label{eq:geometric_limits}
\end{equation}
Substitution in \eqref{eq:QR} gives \eqref{eq:Qinfty}.
\end{proof}

For $\vartheta=(\vartheta_p,\vartheta_q)\in\Torus^2=(\mathbb R/2\pi\mathbb Z)^2$, put
\begin{equation}
u_{a,b}^{(p,q)}=a\log p+b\log q
\label{eq:uab}
\end{equation}
and define
\begin{equation}
K_{R,\vartheta}^{(p,q)}(z)
=
\sum_{(a,b)\in\mathcal S_R}
w_{a,b;R}^{(p,q)}
e^{-i(a\vartheta_p+b\vartheta_q)}
e^{izu_{a,b}^{(p,q)}}.
\label{eq:finite_kernel}
\end{equation}

\section{Limit kernel and pointwise resonance}
\label{sec:limit_kernel}

Define
\begin{equation}
w_{a,b;\infty}^{(p,q)}
=
(1-p^{-1})(1-q^{-1})p^{-a}q^{-b}.
\label{eq:limit_weights}
\end{equation}

\begin{proposition}[Uniform convergence on the real axis]
\label{prop:uniform_convergence}
For every fixed $\vartheta\in\Torus^2$,
\begin{equation}
K_{R,\vartheta}^{(p,q)}(\tau)
\longrightarrow
K_{\infty,\vartheta}^{(p,q)}(\tau)
\label{eq:uniform_kernel_limit}
\end{equation}
uniformly for $\tau\in\mathbb R$, where
\begin{align}
K_{\infty,\vartheta}^{(p,q)}(z)
&=
\sum_{a,b\geq0}
w_{a,b;\infty}^{(p,q)}
e^{-i(a\vartheta_p+b\vartheta_q)}
e^{izu_{a,b}^{(p,q)}}
\label{eq:limit_kernel_series}
\\
&=
\frac{(1-p^{-1})(1-q^{-1})}
{\left(1-p^{-1}e^{i(z\log p-\vartheta_p)}\right)
 \left(1-q^{-1}e^{i(z\log q-\vartheta_q)}\right)}.
\label{eq:limit_kernel_closed}
\end{align}
The series is locally uniformly convergent in the half-plane $\operatorname{Im}z>-1$.
\end{proposition}

\begin{proof}
Extend the finite weight vector by zero outside $\mathcal S_R$.  The exhaustion and the convergence of $W_R^{(p,q)}$ imply
\begin{equation}
\sum_{a,b\geq0}
\left|w_{a,b;R}^{(p,q)}-w_{a,b;\infty}^{(p,q)}\right|
\longrightarrow0.
\label{eq:l1_weight_convergence}
\end{equation}
This proves uniform convergence on the real axis.  The geometric series give \eqref{eq:limit_kernel_closed}; the same argument on compact subsets of $\operatorname{Im}z>-1$ proves local uniform convergence there.
\end{proof}

\begin{proposition}[Pointwise resonance]
\label{prop:pointwise_resonance}
For real $\tau$,
\begin{equation}
\left|K_{\infty,\vartheta}^{(p,q)}(\tau)\right|\leq1.
\label{eq:pointwise_bound}
\end{equation}
Moreover,
\begin{equation}
\sup_{\tau\in\mathbb R}
\left|K_{\infty,\vartheta}^{(p,q)}(\tau)\right|=1.
\label{eq:global_supremum}
\end{equation}
For $\vartheta=(0,0)$ equality occurs at $\tau=0$.
\end{proposition}

\begin{proof}
For $0<r<1$ and real $t$,
\begin{equation}
|1-re^{it}|\geq1-r.
\label{eq:elementary_denominator_bound}
\end{equation}
Applying this to both factors in \eqref{eq:limit_kernel_closed} proves \eqref{eq:pointwise_bound}.  Unique factorization gives
\begin{equation}
m\log p+n\log q=0,\quad (m,n)\in\mathbb Z^2
\quad\Longrightarrow\quad
m=n=0.
\label{eq:log_independence}
\end{equation}
Consequently the line $\tau(\log p,\log q)$ is dense modulo $2\pi$ in $\Torus^2$.  It approaches the prescribed phase $\vartheta$ arbitrarily closely, proving \eqref{eq:global_supremum}.
\end{proof}

\section{Equidistribution of zeta-zero phases}
\label{sec:zero_equidistribution}

Let $\rho=\beta+i\gamma$ run through the nontrivial zeros of $\zeta$, with $0<\gamma\leq T$, counted with multiplicity, and let $N(T)$ be their number.

\begin{theorem}[Joint zero-phase equidistribution]
\label{thm:joint_equidistribution}
For distinct primes $p,q$, the vectors
\begin{equation}
\left(\gamma\log p,\gamma\log q\right)\pmod{2\pi}
\label{eq:zero_phase_vectors}
\end{equation}
are uniformly distributed in $\Torus^2$ as $T\to\infty$.
\end{theorem}

\begin{proof}
Hlawka's theorem states that, for every fixed nonzero real number $\alpha$, the sequence $\alpha\gamma$ is uniformly distributed modulo one; see \cite{FordZaharescu2005}.  For every $(m,n)\in\mathbb Z^2\setminus\{(0,0)\}$, equation \eqref{eq:log_independence} gives
\begin{equation}
\frac{m\log p+n\log q}{2\pi}\neq0.
\label{eq:nonzero_character_frequency}
\end{equation}
Hlawka's theorem applied to every nontrivial character, followed by the multidimensional Weyl criterion, proves the assertion.
\end{proof}

\section{Exact mean-square contraction}
\label{sec:mean_square_contraction}

\begin{theorem}[Finite-shell contraction on the zero ordinates]
\label{thm:finite_shell_contraction}
For every fixed shell $R$ and phase $\vartheta$,
\begin{equation}
\lim_{T\to\infty}
\frac1{N(T)}
\sum_{0<\gamma\leq T}
\left|K_{R,\vartheta}^{(p,q)}(\gamma)\right|^2
=
Q_R^{(p,q)}.
\label{eq:finite_shell_zero_energy}
\end{equation}
Consequently,
\begin{equation}
\lim_{R\to\infty}\lim_{T\to\infty}
\frac1{N(T)}
\sum_{0<\gamma\leq T}
\left|K_{R,\vartheta}^{(p,q)}(\gamma)\right|^2
=
Q_\infty^{(p,q)}<1.
\label{eq:double_limit_contraction}
\end{equation}
\end{theorem}

\begin{proof}
Expanding the square in \eqref{eq:finite_shell_zero_energy} gives diagonal terms
\begin{equation}
\sum_{(a,b)\in\mathcal S_R}
\left(w_{a,b;R}^{(p,q)}\right)^2
=Q_R^{(p,q)}
\label{eq:diagonal_energy}
\end{equation}
and off-diagonal characters of the vector \eqref{eq:zero_phase_vectors}.  Theorem~\ref{thm:joint_equidistribution} makes every off-diagonal average vanish.  Proposition~\ref{prop:coefficient_energy} then gives \eqref{eq:double_limit_contraction}.
\end{proof}

\begin{corollary}[Limit-kernel contraction]
\label{cor:limit_kernel_contraction}
For every fixed phase $\vartheta$,
\begin{equation}
\lim_{T\to\infty}
\frac1{N(T)}
\sum_{0<\gamma\leq T}
\left|K_{\infty,\vartheta}^{(p,q)}(\gamma)\right|^2
=Q_\infty^{(p,q)}.
\label{eq:limit_kernel_zero_energy}
\end{equation}
\end{corollary}

\begin{proof}
The uniform convergence in Proposition~\ref{prop:uniform_convergence} permits passage from the finite kernels to the limit kernel inside the normalized average.
\end{proof}

Define the torus kernel
\begin{equation}
\mathcal K_{p,q}(x,y)
=
\frac{(1-p^{-1})(1-q^{-1})}
{(1-p^{-1}e^{ix})(1-q^{-1}e^{iy})}.
\label{eq:torus_kernel}
\end{equation}

\begin{theorem}[Iteration with one fixed pair]
\label{thm:fixed_pair_iteration}
For every integer $r\geq1$,
\begin{align}
E_r^{(p,q)}
&:=
\lim_{T\to\infty}
\frac1{N(T)}
\sum_{0<\gamma\leq T}
\left|K_{\infty,\vartheta}^{(p,q)}(\gamma)\right|^{2r}
\label{eq:Er_definition}
\\
&=
\frac1{(2\pi)^2}
\int_0^{2\pi}\int_0^{2\pi}
|\mathcal K_{p,q}(x,y)|^{2r}\,\dd x\,\dd y.
\label{eq:Er_torus_integral}
\end{align}
Moreover,
\begin{equation}
E_1^{(p,q)}=Q_\infty^{(p,q)},
\qquad
E_r^{(p,q)}\longrightarrow0
\quad(r\to\infty).
\label{eq:iterated_contraction}
\end{equation}
\end{theorem}

\begin{proof}
Equation \eqref{eq:Er_torus_integral} follows from Theorem~\ref{thm:joint_equidistribution}.  Proposition~\ref{prop:pointwise_resonance} gives $|\mathcal K_{p,q}|\leq1$, with equality only at $(0,0)$ modulo $2\pi$.  This equality set has Haar measure zero.  Dominated convergence proves the second assertion in \eqref{eq:iterated_contraction}.
\end{proof}

\begin{corollary}[The pair $(5,7)$]
\label{cor:five_seven}
For $p=5$ and $q=7$,
\begin{equation}
K_{\infty,0}^{(5,7)}(\tau)
=
\frac{24/35}
{\left(1-5^{-1}e^{i\tau\log5}\right)
 \left(1-7^{-1}e^{i\tau\log7}\right)}
\label{eq:five_seven_kernel}
\end{equation}
and
\begin{equation}
\boxed{
\lim_{T\to\infty}
\frac1{N(T)}
\sum_{0<\gamma\leq T}
\left|K_{\infty,0}^{(5,7)}(\gamma)\right|^2
=\frac12.}
\label{eq:five_seven_half}
\end{equation}
The moments with exponent $2r$ tend to zero in the iterated order displayed in Theorem~\ref{thm:fixed_pair_iteration}.
\end{corollary}

\begin{remark}[One prime versus one dilation]
\label{rem:one_prime}
A single dilation is unitary.  A normalized average over the powers of one prime is a different operator: its limiting mean-square factor on the zero ordinates is $(p-1)/(p+1)<1$.  These two statements must not be identified.
\end{remark}

\section{Mangoldt resonances and the linear bridge}
\label{sec:mangoldt_bridge}

For a nontrivial zero write
\begin{equation}
\delta_\rho=\beta-\frac12.
\label{eq:horizontal_displacement}
\end{equation}
For fixed $T$ with $N(T)>0$, define the existing finite-height exponential difference
\begin{equation}
F_T(y)
=
\frac1{N(T)}
\sum_{0<\gamma\leq T}
\left(e^{(\delta_\rho+i\gamma)y}-e^{i\gamma y}\right).
\label{eq:FT}
\end{equation}

Let $\Lambda(x)$ equal $\log \ell$ when $x=\ell^k$ is an integral prime power and equal zero for all other real $x>1$.

\begin{proposition}[Mangoldt response of a finite shell]
\label{prop:mangoldt_response}
Assume $(0,0)\in\mathcal S_R$.  Write $n_{a,b}=p^aq^b$ and $w_{0,0;R}=w_{0,0;R}^{(p,q)}$.  Then
\begin{align}
&\frac1T
\left(
\sum_{0<\gamma\leq T}
K_{R,\vartheta}^{(p,q)}(\gamma-i\delta_\rho)
-w_{0,0;R}N(T)
\right)
\label{eq:complete_mangoldt_response}
\\
&\qquad\longrightarrow
-\frac1{2\pi}
\sum_{\substack{(a,b)\in\mathcal S_R\\(a,b)\neq(0,0)}}
w_{a,b;R}^{(p,q)}e^{-i(a\vartheta_p+b\vartheta_q)}
\frac{\Lambda(n_{a,b})}{\sqrt{n_{a,b}}},
\notag
\\[1ex]
&\frac1T
\left(
\sum_{0<\gamma\leq T}
K_{R,\vartheta}^{(p,q)}(\gamma)
-w_{0,0;R}N(T)
\right)
\label{eq:ordinate_mangoldt_response}
\\
&\qquad\longrightarrow
-\frac1{2\pi}
\sum_{\substack{(a,b)\in\mathcal S_R\\(a,b)\neq(0,0)}}
w_{a,b;R}^{(p,q)}e^{-i(a\vartheta_p+b\vartheta_q)}
\frac{\Lambda(n_{a,b})}{\sqrt{n_{a,b}}}.
\notag
\end{align}
If $a,b\geq1$, then $\Lambda(p^aq^b)=0$; hence the Mangoldt main term is supported on the two prime-power axes.
\end{proposition}

\begin{proof}
For every fixed $x>1$, Landau's formula gives \cite{Titchmarsh1986}
\begin{equation}
\sum_{0<\gamma\leq T}x^\rho
=
-\frac{T}{2\pi}\Lambda(x)+O_x(\log T).
\label{eq:landau_formula}
\end{equation}
Ford and Zaharescu prove the corresponding fixed-frequency ordinate asymptotic
\cite{FordZaharescu2005}:
\begin{equation}
\frac{1}{T}
\sum_{0<\gamma\leq T}x^{i\gamma}
\longrightarrow
-\frac{1}{2\pi}\frac{\Lambda(x)}{\sqrt{x}},
\qquad x>1 \text{ fixed}.
\label{eq:ford_zaharescu_ordinate_formula}
\end{equation}
Here, as above, $\Lambda(x)$ is zero unless $x$ is an integral prime power.
For completeness, \eqref{eq:ford_zaharescu_ordinate_formula} is the Fourier-coefficient
form of their limiting distribution for the fractional parts of
$\gamma\log x/(2\pi)$; its sign and normalization are therefore fixed by the
coefficient $-\Lambda(x)/(2\pi\sqrt{x})$.

Since
\begin{equation}
K_{R,\vartheta}^{(p,q)}(\gamma-i\delta_\rho)
=
\sum_{(a,b)\in\mathcal S_R}
w_{a,b;R}^{(p,q)}e^{-i(a\vartheta_p+b\vartheta_q)}
n_{a,b}^{\rho-1/2},
\label{eq:complete_kernel_expansion}
\end{equation}
finite summation of \eqref{eq:landau_formula} proves
\eqref{eq:complete_mangoldt_response}.  Likewise,
\begin{equation}
K_{R,\vartheta}^{(p,q)}(\gamma)
=
\sum_{(a,b)\in\mathcal S_R}
w_{a,b;R}^{(p,q)}e^{-i(a\vartheta_p+b\vartheta_q)}
n_{a,b}^{i\gamma},
\label{eq:ordinate_kernel_expansion}
\end{equation}
and finite summation of \eqref{eq:ford_zaharescu_ordinate_formula} proves
\eqref{eq:ordinate_mangoldt_response}.  Unique factorization gives the final assertion.
\end{proof}

\begin{corollary}[Exact linear bridge]
\label{cor:linear_bridge}
For every fixed shell $R$,
\begin{align}
\sum_{(a,b)\in\mathcal S_R}
&w_{a,b;R}^{(p,q)}e^{-i(a\vartheta_p+b\vartheta_q)}
F_T\!\left(u_{a,b}^{(p,q)}\right)
\label{eq:linear_bridge_F}
\\
&=
\frac1{N(T)}
\sum_{0<\gamma\leq T}
\left(
K_{R,\vartheta}^{(p,q)}(\gamma-i\delta_\rho)
-K_{R,\vartheta}^{(p,q)}(\gamma)
\right)
\label{eq:linear_bridge_kernel}
\\
&=o_R\!\left(\frac1{\log T}\right).
\label{eq:linear_bridge_asymptotic}
\end{align}
\end{corollary}

\begin{proof}
The equality of \eqref{eq:linear_bridge_F} and \eqref{eq:linear_bridge_kernel} follows by expanding the definitions.  Proposition~\ref{prop:mangoldt_response} makes the numerator in \eqref{eq:linear_bridge_kernel} equal to $o_R(T)$.  The Riemann--von Mangoldt formula gives $N(T)\asymp T\log T$ \cite{Titchmarsh1986}.
\end{proof}

\section{A positive finite-height energy bridge}
\label{sec:energy_bridge}

Let $m_T(\delta,\tau)$ be the multiplicity of the zero $1/2+\delta+i\tau$ for $0<\tau\leq T$ and $\delta>0$.  For $x>0$, define
\begin{equation}
D_x(\tau)
=
\sum_{\delta>0}
m_T(\delta,\tau)
\left(\cosh(\delta\log x)-1\right).
\label{eq:Dx}
\end{equation}
Functional-equation symmetry pairs $\delta$ with $-\delta$ at the same ordinate.

Let $\dd m(\vartheta)=(2\pi)^{-2}\dd\vartheta_p\dd\vartheta_q$ be normalized Haar measure on $\Torus^2$.  Define
\begin{align}
\mathcal E_{R,T}^{\mathrm{full}}
&=
\frac1{N(T)}
\sum_{0<\gamma\leq T}
\int_{\Torus^2}
\left|K_{R,\vartheta}^{(p,q)}
(\gamma-i\delta_\rho)\right|^2
\dd m(\vartheta),
\label{eq:full_energy}
\\
\mathcal E_{R,T}^{\mathrm{ord}}
&=
\frac1{N(T)}
\sum_{0<\gamma\leq T}
\int_{\Torus^2}
\left|K_{R,\vartheta}^{(p,q)}(\gamma)\right|^2
\dd m(\vartheta),
\label{eq:ordinate_energy}
\\
\Delta_{R,T}^{(p,q)}
&=
\mathcal E_{R,T}^{\mathrm{full}}
-\mathcal E_{R,T}^{\mathrm{ord}}.
\label{eq:energy_difference}
\end{align}

\begin{theorem}[Exact energy bridge]
\label{thm:exact_energy_bridge}
For every finite shell,
\begin{align}
\mathcal E_{R,T}^{\mathrm{ord}}
&=Q_R^{(p,q)},
\label{eq:exact_ordinate_energy}
\\
\Delta_{R,T}^{(p,q)}
&=
\frac{2}{N(T)}
\sum_{0<\tau\leq T}
\sum_{\delta>0}m_T(\delta,\tau)
\sum_{(a,b)\in\mathcal S_R}
\left(w_{a,b;R}^{(p,q)}\right)^2
\left(\cosh(2\delta u_{a,b}^{(p,q)})-1\right)
\label{eq:positive_energy_bridge}
\\
&=
\frac{2}{N(T)}
\sum_{(a,b)\in\mathcal S_R}
\left(w_{a,b;R}^{(p,q)}\right)^2
\sum_{0<\tau\leq T}
D_{\exp(2u_{a,b}^{(p,q)})}(\tau).
\label{eq:energy_bridge_Dx}
\end{align}
If $\mathcal S_R$ contains an index $(a,b)\neq(0,0)$, then
\begin{equation}
\boxed{
\Delta_{R,T}^{(p,q)}=0
\quad\Longleftrightarrow\quad
\operatorname{Re}\rho=\frac12
\text{ for every }0<\operatorname{Im}\rho\leq T.}
\label{eq:energy_bridge_equivalence}
\end{equation}
\end{theorem}

\begin{proof}
Orthogonality of the phase characters gives, for every real $\delta$,
\begin{equation}
\int_{\Torus^2}
\left|K_{R,\vartheta}^{(p,q)}(\gamma-i\delta)\right|^2
\dd m(\vartheta)
=
\sum_{(a,b)\in\mathcal S_R}
\left(w_{a,b;R}^{(p,q)}\right)^2
e^{2\delta u_{a,b}^{(p,q)}}.
\label{eq:phase_orthogonality_energy}
\end{equation}
Setting $\delta=0$ proves \eqref{eq:exact_ordinate_energy}.  Pairing $\delta$ and $-\delta$ proves \eqref{eq:positive_energy_bridge}; Definition~\eqref{eq:Dx} gives \eqref{eq:energy_bridge_Dx}.  Every summand in \eqref{eq:positive_energy_bridge} is nonnegative.  A nonzero displacement makes every term with $u_{a,b}^{(p,q)}>0$ strictly positive, proving \eqref{eq:energy_bridge_equivalence}.
\end{proof}

Define the finite-height curvature
\begin{equation}
F_T''(0)
=
\frac{2}{N(T)}
\sum_{0<\tau\leq T}
\sum_{\delta>0}m_T(\delta,\tau)\delta^2.
\label{eq:curvature}
\end{equation}

\begin{corollary}[Coercive curvature bound]
\label{cor:coercive_curvature}
Let
\begin{equation}
M_{2,R}^{(p,q)}
=
\sum_{(a,b)\in\mathcal S_R}
\left(w_{a,b;R}^{(p,q)}\right)^2
\left(u_{a,b}^{(p,q)}\right)^2.
\label{eq:M2R}
\end{equation}
Then
\begin{equation}
\boxed{
\Delta_{R,T}^{(p,q)}
\geq
2M_{2,R}^{(p,q)}F_T''(0).}
\label{eq:coercive_bridge}
\end{equation}
\end{corollary}

\begin{proof}
Apply $\cosh t-1\geq t^2/2$ termwise in \eqref{eq:positive_energy_bridge}.
\end{proof}

\begin{corollary}[Explicit $(5,7)$ bridge]
\label{cor:explicit_five_seven_bridge}
For the limit weights
\begin{equation}
w_{a,b;\infty}^{(5,7)}=\frac{24}{35}5^{-a}7^{-b},
\label{eq:five_seven_limit_weights}
\end{equation}
set
\begin{equation}
J_{5,7}(\delta)
=
\sum_{a,b\geq0}
\left(w_{a,b;\infty}^{(5,7)}\right)^2
e^{2\delta(a\log5+b\log7)}
=
\frac{(24/35)^2}
{(1-5^{2\delta-2})(1-7^{2\delta-2})}.
\label{eq:J_five_seven}
\end{equation}
For $|\delta|<1/2$ this is finite, and $J_{5,7}(0)=1/2$.  The limiting energy difference satisfies
\begin{equation}
\Delta_{\infty,T}^{(5,7)}
=
\frac1{N(T)}
\sum_{0<\gamma\leq T}
\left(J_{5,7}(\delta_\rho)-\frac12\right)
\label{eq:explicit_energy_difference}
\end{equation}
and
\begin{equation}
\boxed{
\Delta_{\infty,T}^{(5,7)}
\geq
0.204533768513927\ldots\,F_T''(0).}
\label{eq:explicit_coercive_constant}
\end{equation}
More exactly,
\begin{equation}
M_{2,\infty}^{(5,7)}
=
\frac{13}{576}(\log5)^2
+\frac{25}{2304}(\log7)^2
+\frac{1}{1152}\log5\log7
=0.102266884256963\ldots.
\label{eq:explicit_M2}
\end{equation}
\end{corollary}

\begin{proof}
Equation \eqref{eq:J_five_seven} is the product of two geometric series.
Passing to the limit shell in \eqref{eq:phase_orthogonality_energy} gives
\eqref{eq:explicit_energy_difference}.  Direct differentiation of the geometric series, or the identities
\begin{equation}
\sum_{k\geq0}r^k=\frac1{1-r},
\quad
\sum_{k\geq0}kr^k=\frac{r}{(1-r)^2},
\quad
\sum_{k\geq0}k^2r^k=\frac{r(1+r)}{(1-r)^3},
\label{eq:geometric_moments}
\end{equation}
gives \eqref{eq:explicit_M2}.  Corollary~\ref{cor:coercive_curvature} then gives \eqref{eq:explicit_coercive_constant}.
\end{proof}

\section{What the contraction determines}
\label{sec:scope}

\begin{theorem}[Finite horizontal perturbations preserve the asymptotic inputs]
\label{thm:finite_perturbation}
Let a symmetric zero multiset satisfy the Riemann--von Mangoldt count, Hlawka equidistribution, the Landau formula \eqref{eq:landau_formula}, and the mean-square contraction \eqref{eq:five_seven_half}.  Adjoin, at one fixed positive ordinate $\tau$, the symmetric pair
\begin{equation}
\frac12+\delta+i\tau,
\qquad
\frac12-\delta+i\tau,
\qquad 0<\delta<\frac12,
\label{eq:inserted_pair}
\end{equation}
together with the conjugate pair at negative ordinate.  Then all four asymptotic statements remain valid, while for every $T\geq\tau$,
\begin{equation}
F_T''(0)>0,
\qquad
\Delta_{\infty,T}^{(5,7)}>0.
\label{eq:positive_after_insertion}
\end{equation}
Consequently, the mean-square contraction and the stated Mangoldt asymptotics do not imply the energy equality $\Delta_{\infty,T}^{(5,7)}=0$.
\end{theorem}

\begin{proof}
The insertion changes $N(T)$ by at most two and changes every fixed normalized zero-phase average by $O(1/N(T))$.  It therefore preserves the Riemann--von Mangoldt and Hlawka asymptotics and the limit in \eqref{eq:five_seven_half}.  For every fixed $x>1$, its contribution to the Landau sum is
\begin{equation}
x^{1/2+\delta+i\tau}+x^{1/2-\delta+i\tau}
=
2x^{1/2}e^{i\tau\log x}\cosh(\delta\log x),
\label{eq:finite_landau_perturbation}
\end{equation}
which is $O_x(1)$ and is absorbed by the error $O_x(\log T)$ in \eqref{eq:landau_formula}.  The strict inequalities in \eqref{eq:positive_after_insertion} follow from \eqref{eq:curvature} and Theorem~\ref{thm:exact_energy_bridge}.
\end{proof}

\begin{proposition}[Energy transfer preserves the horizontal excess]
\label{prop:energy_transfer}
For arbitrary complex branch amplitudes $z_p,z_q$,
\begin{equation}
\left|\frac{z_p+z_q}{\sqrt2}\right|^2
+
\left|\frac{z_p-z_q}{\sqrt2}\right|^2
=
|z_p|^2+|z_q|^2.
\label{eq:parallelogram_energy}
\end{equation}
Thus contraction of the averaged branch transfers the missing energy to the complementary branch.  Applied to the complete zero branches, the conserved total is $Q_R^{(p,q)}+\Delta_{R,T}^{(p,q)}$, not $Q_R^{(p,q)}$ when a horizontal displacement is present.
\end{proposition}

\begin{proof}
Equation \eqref{eq:parallelogram_energy} is the parallelogram identity.  The second assertion follows from the definition \eqref{eq:energy_difference}.
\end{proof}

\section{Conclusion}
\label{sec:conclusion}

The prime pair $(5,7)$ has the exact zero-ordinate mean-square contraction factor $1/2$, and the moments of every fixed iterate are given by the torus integral \eqref{eq:Er_torus_integral}.  The isolated pointwise resonances do not obstruct this averaged contraction.  The Mangoldt response is supported on the two prime-power axes; the mixed branches $5^a7^b$, $a,b\geq1$, have zero Mangoldt main term.

The linear identity \eqref{eq:linear_bridge_kernel} connects the kernel directly to the existing finite-height exponential difference $F_T$.  Phase orthogonality gives the stronger positive energy identity \eqref{eq:positive_energy_bridge}.  For $(5,7)$ it yields the explicit coercive estimate \eqref{eq:explicit_coercive_constant}, and vanishing of the energy difference is equivalent to critical-line containment up to height $T$.

Theorem~\ref{thm:finite_perturbation} determines the logical relation between these statements: the zero-ordinate contraction and the stated Landau--Mangoldt asymptotics are unchanged by a finite symmetric horizontal perturbation, whereas the positive energy bridge responds immediately.  Energy conservation transfers the corresponding excess between complementary branches but does not remove it.

\begin{thebibliography}{99}

\bibitem{FordZaharescu2005}
K.~Ford and A.~Zaharescu,
On the distribution of imaginary parts of zeros of the Riemann zeta function,
\emph{J. Reine Angew. Math.} \textbf{579} (2005), 145--158,
\href{https://doi.org/10.1515/crll.2005.2005.579.145}{doi:10.1515/crll.2005.2005.579.145}; arXiv:math/0405459.

\bibitem{Titchmarsh1986}
E.~C.~Titchmarsh,
\emph{The Theory of the Riemann Zeta-function},
2nd ed., revised by D.~R.~Heath-Brown,
Oxford University Press, 1986.

\end{thebibliography}

\end{document}
